\documentclass[10pt,twocolumn]{article}
\usepackage[letterpaper,margin=1in]{geometry}
\usepackage[T1]{fontenc}
\usepackage[utf8]{inputenc}
\usepackage{lmodern}
\usepackage[protrusion=true,expansion=false]{microtype}
\usepackage{xcolor}
\usepackage{amsmath}
\usepackage{amssymb}
\usepackage{array}
\usepackage{balance}
\usepackage{booktabs}
\usepackage{enumitem}
\usepackage[pdfencoding=pdfdoc]{hyperref}
\usepackage[numbers,sort&compress]{natbib}

\definecolor{linkblue}{RGB}{0,42,120}
\hypersetup{
  colorlinks=true,
  linkcolor=linkblue,
  citecolor=linkblue,
  urlcolor=linkblue
}

\AtBeginDocument{\fontsize{9.55pt}{10.95pt}\selectfont}
\newcolumntype{L}[1]{>{\raggedright\arraybackslash}p{#1}}
\geometry{letterpaper,top=0.72in,bottom=0.82in,left=0.7in,right=0.7in}
\setlength{\columnsep}{0.24in}
\setlength{\parskip}{1.2pt}
\setlength{\parindent}{1em}
\setlength{\abovedisplayskip}{5pt plus 1pt minus 1pt}
\setlength{\belowdisplayskip}{5pt plus 1pt minus 1pt}
\setlength{\abovedisplayshortskip}{3pt plus 1pt minus 1pt}
\setlength{\belowdisplayshortskip}{4pt plus 1pt minus 1pt}
\setlength{\tabcolsep}{4pt}
\renewcommand{\arraystretch}{1.06}
\setlist[itemize]{leftmargin=1.15em,itemsep=1pt,topsep=2.5pt,parsep=0pt,partopsep=0pt}
\setlist[enumerate]{leftmargin=1.35em,itemsep=1pt,topsep=2.5pt,parsep=0pt,partopsep=0pt}
\newlist{facts}{description}{1}
\setlist[facts]{
  style=nextline,
  leftmargin=0pt,
  labelwidth=0pt,
  labelsep=0pt,
  font=\normalfont\bfseries,
  itemsep=1.5pt,
  topsep=2pt,
  parsep=0pt,
  partopsep=0pt
}
\setlength{\emergencystretch}{2em}
\clubpenalty=10000
\widowpenalty=10000
\displaywidowpenalty=10000
\renewcommand{\bibfont}{\footnotesize}
\setlength{\bibsep}{1pt plus 0.2ex}
\raggedbottom

\newcommand{\policystrata}{\mbox{\textsc{PolicyStrata}}}
\newcommand{\bench}{\textsc{\mbox{DataPolicy}\allowbreak\mbox{DriftBench}}}


\title{\policystrata: Responsibility-Scoped Testing for Cross-Layer Policy Drift in LLM Data Agents}
\author{Zachary Roth\\Raintree Technology\\\texttt{zacharyroth@pm.me}}
\date{June 2026; revised July 2026}
\hypersetup{
  pdftitle={PolicyStrata: Responsibility-Scoped Testing for Cross-Layer Policy Drift in LLM Data Agents},
  pdfauthor={Zachary Roth},
  pdfsubject={Responsibility-scoped testing for cross-layer policy drift in LLM data agents},
  pdfkeywords={PolicyStrata, LLM data agents, policy drift, regression testing}
}

\begin{document}
\twocolumn[
\begin{@twocolumnfalse}
\maketitle
\vspace{-1.5em}
\begin{abstract}
LLM data agents copy policy obligations across a model-visible manifest, grammar, validator, query compiler, database controls, and result-release logic. Each component can be locally correct while a version or semantic mismatch between components exposes the wrong capability, lowers an authorized plan into unauthorized SQL, or releases a disallowed result.

We present \policystrata, a deterministic regression tester for this cross-layer policy drift. A structured policy and surface contract define which obligations each transition must preserve. \policystrata{} replays or generates a principal, semantic query, database state, and surface-version vector; checks the transitions in execution order; and emits a replayable witness for the first violated contract. Its v1 benchmark contains 22 mutation operators over three synthetic analytics domains. The 1720/1720 injected-fault result is a construction-consistency check over that taxonomy, not production recall. More discriminating results are that a specification-derived conventional test suite misses 141/1720 cases, a stacked set of common point controls misses 159/1720, and a spec-blind suite agrees with \policystrata{} on 39/42 cases. Measured against an independently authored data-agent vulnerability taxonomy, the operator registry covers two of eight classes and one partially, and 401 of its own cases describe update-induced drift that an attack taxonomy has no reason to name. Static retargeting to four independently maintained open-source stacks measured one content-level false positive among 74 real-SQL traces and exposed five adapter gaps, all since fixed. Executing all 20 row-level-security policies in one target's frozen migrations passed 13 live containment checks, and weakening a single real predicate failed exactly the checks covering it. A source-frozen MetricFlow run reproduced 68 upstream-authored expected-SQL cases, while a deployment-linked BetterOff pilot bound the adapter to the exact production revision and passed 33 live denial-boundary probes; three authenticated probes remained unavailable. Source-contract replay reproduced three historical BetterOff pre-fix revisions, two mapping to v1 and one outside it. These results establish reproducible fault-model coverage and a concrete retargeting path; they do not establish field recall or independently operated deployment effectiveness.
\end{abstract}
\vspace{0.8em}
\end{@twocolumnfalse}
]

\section{Introduction}

An LLM data agent turns a request into a structured plan, executable query, and released answer. Policy enters this path more than once: a capability manifest tells the model what it may request; a grammar defines valid plans; a validator accepts or rejects a plan; a compiler turns it into SQL; database policy constrains execution; and a release check decides what may leave the trusted boundary. We call these six components the \emph{policy pipeline}. They use different representations and are often deployed on different schedules.

The resulting failure is not ordinary component failure. Every component may pass its own tests while the transition between two components drops tenant identity, changes a metric definition, uses a stale alias, loses lineage, or releases a result under an obsolete rule. We call a mismatch between a canonical obligation and a policy-pipeline surface \emph{divergence}; when an update creates it, we call it \emph{policy drift}.

\paragraph{Worked example.}
Suppose analyst Alice belongs to tenant $A$ and may request tenant-$A$ ticket aggregates, but not tenant-$B$ rows or raw customer email. The canonical policy is version 7. The validator also runs v7, while the compiler still runs v5 and uses \texttt{legacy\_tenant\_id}. Alice asks for ``escalations by region.'' The validator accepts the semantic plan
\[
\langle\texttt{escalated\_tickets},[\texttt{region}],
\texttt{last\_month},100\rangle .
\]
Compiling that accepted plan is a \emph{lowering}: a translation from the typed plan into SQL. The stale compiler binds the tenant obligation to the wrong column. On a two-row database containing one distinguishing tenant-$B$ escalation, the emitted SQL reads both tenants. Database row-level security (RLS) may block the extra row, but that containment does not make the compiler correct.

\policystrata{} records the six deployed versions as a \emph{version vector}, here
\[
\begin{aligned}
\vec v={}&(v_{\rm manifest},v_{\rm grammar},v_{\rm validator},\\
          &v_{\rm compiler},v_{\rm db},v_{\rm release})
          =(7,7,7,5,7,7).
\end{aligned}
\]
It evaluates the canonical plan and the lowered query, checks each transition in order, and reports the compiler as the first violated contract. The witness retains Alice, the plan, the version vector, the distinguishing rows, the SQL, the RLS outcome, and the release decision. This example separates three facts that a final-answer assertion would collapse: the compiler violated tenant preservation, RLS contained the violation, and no external disclosure occurred.

The right oracle is therefore not literal equality between layers. A manifest may intentionally expose less than a validator accepts, and RLS may intentionally be stricter than an application check. \policystrata{} assigns each surface a responsibility and checks only the obligations that surface accepts or emits.

\paragraph{Scope and evidence boundary.}
\policystrata{} is a regression tester and scanner, not an authorization boundary. It assumes a human-reviewed canonical policy and trusted adapters that map implementation artifacts into the checking model. The v1 benchmark covers single-request, read-only analytics. Its deterministic generator selects plans from the same public operator taxonomy used by the simulator, so its perfect kill count measures internal coverage. A maintainer-operated deployment study verifies source-to-deployment binding and selected live denial behavior, but publishes only aggregate results, does not read customer rows, and lacks an authenticated synthetic principal. No evaluation observes an independently operated deployment or estimates recall on unknown production incidents.

\paragraph{Claims and non-claims.}
\begin{center}
\small
\begin{tabular}{@{}L{0.43\columnwidth}L{0.47\columnwidth}@{}}
\toprule
Claim & Non-claim \\
\midrule
Responsibility contracts detect and localize faults represented by the implemented model & The contracts do not prove a deployment's canonical policy is correct \\
The artifact covers its 22 operators and distinguishes gaps missed by two spec-derived comparators & 1720/1720 is not recall on unknown faults \\
The scanner can target real artifacts and an exact deployed revision & Deployment binding and denial probes are not authenticated cross-tenant validation \\
Every reported witness has a modeled contract violation in the checked state space & This is not a mechanized soundness proof or global completeness result \\
\bottomrule
\end{tabular}
\end{center}

\paragraph{Contributions.}
This paper contributes:

\begin{itemize}
  \item a model of a data agent as six non-equivalent policy surfaces with explicit exposure, validation, lowering, containment, and release contracts;
  \item a deterministic generator and replay checker that produces first-transition witnesses, plus counterfactual repair checks for attribution and bounded replay reduction for witness size;
  \item \bench, with 22 named operators, explicit suite provenance, equivalent/invalid accounting, freeze manifests, two spec-derived comparators, and separate synthetic, spec-blind, public-fault, brownfield, and historical-revision evidence;
  \item a documented retargeting contract plus an aggregate, deployment-linked read-only study whose private operational identifiers are withheld.
\end{itemize}

\section{Background and Related Work}

Six research lines touch parts of the policy pipeline, but they answer different questions.

\paragraph{Access-control analysis and mutation testing.}
Margrave computes semantic differences between two access-control policies and can explain requests whose decisions change \citep{margrave}. XACML mutation work defines policy fault operators, test adequacy, and equivalent-mutant handling \citep{xacml_mutation}. Cedar compares an authorization implementation with an independently specified model through formal and differential testing \citep{cedar_vgd}. These methods can analyze a policy language more deeply than \policystrata{} and may prove or systematically test decision equivalence inside that language. They do not by themselves connect a model-visible metric alias, semantic plan, SQL lowering, runtime database role, and result-release decision. \policystrata{} borrows mutation adequacy and differential-oracle ideas, but makes the unit of analysis a heterogeneous transition in the policy pipeline.

\paragraph{Translation validation and property-preserving compilation.}
\policystrata{}'s compiler contract is an instance of translation validation: rather than proving a compiler correct once, check each individual compilation against its input \citep{translation_validation}. Secure-compilation work sharpens what such a check should preserve, distinguishing full abstraction from the weaker robust preservation of specific property classes across a compilation boundary \citep{secure_compilation}. Authorization-preserving lowering is deliberately the weaker, property-class form: we do not ask a compiler to preserve all observational equivalences, only the declared obligations of principal, tenant, purpose, policy version, lineage, and metric semantics. What \policystrata{} adds is neither a new preservation theorem nor a proof technique, but the observation that a data agent runs several such lowerings between representations that are separately versioned, so a preservation failure must be attributed to one transition rather than to the pipeline as a whole.

\paragraph{SQL checking and generated databases.}
Distilled test-suite accuracy constructs databases that distinguish candidate SQL from a reference query \citep{test_suite_accuracy}; VeriEQL and SpotIt+ reason about bounded query equivalence under supported SQL semantics and constraints \citep{verieql,spotit}. Such tools give stronger SQL-level semantic evidence than finite execution alone. They do not decide whether a query was exposed to the right principal, retained tenant identity and policy version, used authorized lineage, or crossed an allowed release boundary. Conversely, \policystrata{} does not prove general SQL equivalence: its database states are finite witnesses, and a supported SQL verifier could replace that one checker.

\paragraph{Database and agent enforcement.}
Beacon checks consistency within DBMS access-control metadata and behavior \citep{beacon}, and row-level security is the containment mechanism our database contract observes \citep{postgresql_rls}. Agent enforcement systems place deterministic policy or privilege controls around tool execution \citep{pcas,progent}; Cordon adds a transactional boundary that stages irreversible effects until a composed execution flow validates \citep{cordon}; C-Trace compiles a subset of GDPR into predicates over execution traces and rejects non-compliant actions at runtime \citep{c_trace}; Data Flow Control pushes tuple-level release policy into the DBMS through provenance-aware query rewriting \citep{data_flow_control}; and role-conditioned refusal work evaluates model-level compliance \citep{role_conditioned_refusals}. These systems can block behavior at runtime. \policystrata{} instead tests whether separately maintained enforcement surfaces still agree with their assigned contracts, and it reports an upstream compiler violation even when RLS safely contains it. The relationship is complementary in both directions: an enforcement layer is a surface \policystrata{} can observe, and \policystrata{} is a way to detect that such a layer has drifted from the specification it was deployed to enforce.

\paragraph{Agent conformance testing and drift.}
The closest recent work shares our premise that agent safety degrades between components rather than inside one. Constraint drift argues that safety-critical constraints lose force as they pass through memory, delegation, communication, tool use, and audit, and that asserting a constraint once is not the same as maintaining it \citep{constraint_drift}. \policystrata{} is a narrow, executable instance of that position for a single-agent data pipeline: six named surfaces, declared obligations per transition, and a checker that reports where an obligation stopped holding. AgentRFC extracts normative clauses from agent protocol specifications into a typed IR, model-checks them, and replays counterexamples against live SDKs \citep{agentrfc}; its extract-check-replay shape is close to ours, but its subject is protocol conformance between agents rather than policy obligations between representations inside one agent. TDAD treats prompts as compiled artifacts and applies hidden test splits and semantic mutation testing to agent specifications \citep{tdad}, which is the same adequacy argument we make for a cross-layer operator registry. AgentRaft detects data over-exposure by building a cross-tool call graph, tracking taint at runtime, and adjudicating with a model committee \citep{agentraft}; it covers tool-to-tool flows that \policystrata{} does not model, while \policystrata{} stays deterministic, needs no model at scoring time, and localizes to a versioned transition rather than to a flow.

\paragraph{Generation and task benchmarks.}
Spider, BIRD, DAB, BEAVER, and DAComp measure SQL or data-task capability \citep{spider,bird,dab,beaver,dacomp}, while $\tau$-bench and ST-WebAgentBench measure whether agents follow domain rules and safety policies during interaction \citep{taubench,stwebagentbench}. PICARD and structured decoders reduce malformed outputs \citep{picard,outlines,xgrammar,jsonschemabench}. Correct, rule-following, or grammar-valid output can still select the wrong tenant, implement the wrong business metric, or violate release policy once it is lowered to SQL. \policystrata{} therefore starts at typed semantic IR; stochastic natural-language-to-plan reachability is reported separately and is not part of its deterministic score. A recent security study of data agents supplies the independent fault vocabulary we test our registry against \citep{data_agents_under_attack}; LASM supplies a second, structural taxonomy derived from 116 agent-security papers \citep{lasm}.

\begin{center}
\scriptsize
\setlength{\tabcolsep}{2pt}
\begin{tabular}{@{}L{0.24\columnwidth}L{0.31\columnwidth}L{0.35\columnwidth}@{}}
\toprule
System or method & Evidence it can provide & Boundary relative to \policystrata{} \\
\midrule
Margrave & semantic policy-version difference & one access-control language; no SQL or release path \\
XACML mutation & test adequacy against policy faults & mutates one policy representation \\
Cedar VGD & model/engine differential evidence & authorization decisions, not business semantics \\
Translation validation & per-compilation correctness & one lowering; no versioned surface attribution \\
Distilled SQL suites & distinguishing database states & query results, not authority or lineage \\
SQL verifiers & bounded equivalence evidence & supported SQL fragment only \\
Beacon & DBMS catalog/behavior inconsistency & database internals, not upstream agent surfaces \\
Runtime monitors & enforcement on observed actions & no first-transition diagnosis by itself \\
AgentRFC & protocol conformance via IR and replay & between agents, not between pipeline representations \\
AgentRaft & tool-flow over-exposure detection & model-adjudicated flows, not versioned lowerings \\
\policystrata{} & cross-representation witness and attribution & contract- and taxonomy-relative; no global proof \\
\bottomrule
\end{tabular}
\end{center}

The table is a capability comparison, not an empirical ranking: these systems were not run as drop-in competitors because they accept different languages and prove or test different properties. A fair composition would use Margrave or Cedar as the authorization oracle, VeriEQL as the compiler-semantic checker, and database/runtime monitors as observed surfaces inside \policystrata{}'s transition harness. The novelty boundary is therefore narrow. \policystrata{} is not a new policy logic, SQL verifier, or reference monitor. It composes evidence across versioned, non-equivalent representations and assigns a failed obligation to the first transition responsible for preserving it.

\section{Model and Guarantee Boundary}

Let a data-agent pipeline have runtime objects:

\[
X_0 \to X_1 \to X_2 \to X_3 \to X_4
\]

The representations may correspond to a model-visible request, a typed semantic plan, a validated executable plan, a SQL query and execution context, and a released analytical result. The surfaces that handle those objects are related but not interchangeable:

\begin{center}
\vspace{-0.4em}
\footnotesize
\renewcommand{\arraystretch}{0.92}
\begin{tabular}{@{}L{0.92\columnwidth}@{}}
\toprule
Request or latent intent \\
$\downarrow$ \\
\textbf{Manifest or grammar}: expose only authorized capabilities \\
$\downarrow$ \\
\textbf{Semantic-plan validator}: accept supported authorized plans and reject unauthorized plans \\
$\downarrow$ \\
\textbf{Compiler or lowering}: preserve tenant, purpose, policy version, lineage, and semantics \\
$\downarrow$ \\
\textbf{Database policy and RLS}: contain unauthorized executable operations \\
$\downarrow$ \\
\textbf{Release policy}: release only allowed result-lineage pairs \\
\bottomrule
\end{tabular}
\vspace{-0.7em}
\end{center}

The surface-specific contracts used by \policystrata{} are summarized in Table 1. The formal definitions appear in Appendix~\ref{app:formalization}.

\begin{center}
\begin{minipage}{\columnwidth}
\centering
\footnotesize
\renewcommand{\arraystretch}{0.96}
\textbf{Table 1: Policy-bearing surfaces and their responsibility-scoped obligations.}\\[2pt]
\begin{tabular}{@{}L{0.30\columnwidth}L{0.60\columnwidth}@{}}
\toprule
Surface & Responsibility \\
\midrule
Manifest or grammar & Do not expose capabilities whose reachable operations are unauthorized; underexposure may be intentional. \\
Validator & Accept authorized objects inside the declared support envelope and reject unauthorized objects. \\
Compiler & Preserve principal, tenant, purpose, policy version, lineage, and business semantics during lowering. \\
Database policy & Contain unauthorized executable operations under the runtime database role. \\
Release policy & Release only result-lineage pairs allowed to leave the trusted boundary. \\
\bottomrule
\end{tabular}
\end{minipage}
\end{center}

A grammar is a language recognizer. A validator is a decision procedure. A compiler is a partial lowering function. RLS is an execution-time reference monitor. A manifest is a projection shown to the model. Treating these as one kind of object obscures the bugs we want to find.

At version $t$, the canonical policy separates authorization, semantic, and release specifications: $P^t=(P_{\mathsf{auth}}^t,P_{\mathsf{sem}}^t,P_{\mathsf{rel}}^t)$. The authorization specification defines who may perform operations over rows, columns, metrics, tenants, and purposes. The semantic specification defines the meanings of metrics, dimensions, aliases, grains, joins, filters, and time windows. The release specification defines whether a result and its lineage may leave a trusted boundary. These are intentionally separate: a correct RBAC decision cannot prove that ``net revenue'' used the right denominator, and a semantically correct aggregate may still be unsafe to disclose.

\paragraph{Exposure soundness.}
A model-visible manifest or grammar is exposure-sound if every operation reachable through an exposed capability is authorized for the principal and context. This does not require every authorized capability to be exposed.

\paragraph{Declared completeness.}
Over-restrictive behavior is a defect only inside the advertised support envelope. This avoids treating intentional underexposure, feature flags, stricter release suppression, or unsupported-but-authorized requests as bugs.

\paragraph{Authorization-preserving lowering.}
Lowering must also preserve security-relevant context: principal, tenant, purpose, jurisdiction, policy version, release boundary, and time semantics cannot be silently dropped. This property catches compilers that accept a valid plan but emit SQL that reads the wrong tenant, expands a metric into forbidden columns, or uses stale identity keys.

\paragraph{Semantic translation validation.}
\policystrata{} searches for admissible database states where the lowered query and the reference semantic plan produce different observations. The comparison must define bag versus set semantics, ordering, floating-point tolerance, NULL and three-valued logic, errors, timestamp and timezone behavior, and nondeterministic functions. Testing finitely many generated states finds distinguishing witnesses; it does not prove general SQL equivalence unless paired with a verifier such as bounded SQL-equivalence checking \citep{verieql,spotit}. Semantic equality is not sufficient for authorization: two queries can return the same visible result while one scans forbidden rows or columns. \policystrata{} therefore checks both denotational behavior and dependency or lineage conformance.

\paragraph{Release conformance.}
For v1, release is stateless and single-query: a result and its lineage may leave a boundary only if the release policy permits that pair for the principal and context. A release boundary may be the human user, the LLM context, an external model provider, a log, a cache, or another tool. Raw unauthorized rows entering an untrusted model context can be a disclosure even if the final user-facing answer is suppressed.

\paragraph{Witness classes.}
\policystrata{} reports over-permissive exposure, over-restrictive rejection, lowering violation, semantic drift, and unsafe release. The v1 mutation registry instantiates four of these classes; over-restrictive behavior is supported by the detector but absent from the 22-operator generated benchmark.

\paragraph{Drift witness.}
A drift witness records the principal, request, semantic IR, version vector, lowered SQL, observed database result, release decision, surface responsibilities, contract decisions, first violated transition, and explanatory reasons. The current emitted witness does not reduce arbitrary source code, policy clauses, database rows, or SQL structure. Its bounded reducer only removes semantic-IR dimensions and filters or resets a non-default limit, replaying after each edit to preserve witness class, localized surface, containment layer, release decision, and relevant semantic or database evidence.

\paragraph{Divergence, drift, and version skew.}
A static mismatch present from initial deployment is policy divergence. Update-induced drift exists when the pipeline satisfied expected contracts before an update but fails afterward. A version vector is simply the ordered version identifier for each of the six surfaces. It lets a test reproduce partial rollouts and stale caches instead of recording only one application version. \policystrata{} can detect divergence and skew only when the canonical specification or another independent oracle disagrees with the implementation. It cannot detect a common-mode error where every surface faithfully implements a bad canonical policy.

\paragraph{Contract-relative soundness.}
For an emitted trace $T$, let $\mathsf{Witness}(T)$ mean the detector reports a non-clean class and let $\mathsf{Violation}(T)$ mean at least one modeled surface contract fails or an unauthorized result is released. The checked invariant is
\[
\mathsf{Witness}(T)\Rightarrow\mathsf{Violation}(T).
\]
The implementation tests this property with 400 property-generated examples and an exhaustive sweep of applicable operator/domain combinations over 25 seeds. No counterexample was found. This is evidence that the detector does not invent a witness relative to its own contracts and simulator. It is not a mechanized proof, does not validate the canonical policy or adapters, and does not imply the converse for faults outside the 22 operators.

\section{Checking Procedure}

\policystrata{} separates the reference model, observed surfaces, and labels. The canonical policy is YAML validated into principals, roles, metrics, dimensions, aliases, grains, cost and row bounds, and tenant scopes. A second YAML document declares the six surface versions, each surface's responsibility, and the obligations passed between surfaces. A semantic query is a typed tuple of metric, dimensions, filters, time range, grain, and limit. Static suites or imported JSONL traces supply the remaining test inputs.

\paragraph{Specification language.}
The language is deliberately a typed data schema rather than a new policy calculus. Authored policy and surface files reject unknown fields. A shortened, valid fragment is:

\begin{center}
\begin{minipage}{0.94\columnwidth}
\footnotesize
\begin{verbatim}
roles:
  analyst:
    allowed_metrics: [ticket_count]
    allowed_dimensions: [region]
    max_rows: 1000
contracts:
  compiler:
    mode: sql_lowering
    accepts_obligations: [tenant_scope]
    emits_obligations: [sql_semantics]
\end{verbatim}
\end{minipage}
\end{center}

Metric declarations bind a semantic name to an expression, source table, columns, aliases, roles, grain, and estimated cost. Principal declarations bind a role and tenant identifiers. The surface file does not encode executable enforcement; it declares what an adapter must observe and what the checker will hold that transition responsible for. Imported JSONL then supplies an observed principal, SQL, optional semantic IR, tenant bindings, release outcome, source, and regression label. Missing semantic IR disables semantic-oracle checks for that trace rather than synthesizing a meaning.

\paragraph{Reference and implementation paths.}
The reference interpreter decides whether the semantic query is authorized and computes the intended metric and release conditions. The implementation path simulates or imports manifest exposure, grammar membership, validator output, compiled SQL, database effects, and release output. The reference interpreter does not call the SQL compiler. This separation prevents compiler behavior from defining its own oracle, but a shared error in the policy file or adapters remains possible.

\paragraph{Cross-layer conformance analysis.}
For each case, the checker performs the following steps:

\begin{enumerate}
  \item Load principal $p$, query $q$, database state $D$, canonical policy version $t$, and surface-version vector $\vec v$.
  \item Evaluate $q$ with the reference authorization and semantic interpreters.
  \item Replay the six surfaces in execution order. Each adapter emits its decision and the obligations it passes downstream.
  \item Evaluate only the declared contract at each surface. For example, the compiler check compares tenant binding, metric expression, grain, time semantics, cost, and lineage; it does not require its representation to equal the validator's.
  \item Select the first failing contract in the fixed order manifest, grammar, validator, compiler, database, release. Later enforcement can be recorded as containment without erasing the earlier violation.
  \item Replay bounded reductions and emit the smallest witness reached under the reducer's move set.
\end{enumerate}

The selection rule is deterministic. If $F_i(T)$ is the Boolean contract result for surface $i$, localization returns $\min\{i\mid \neg F_i(T)\}$ in pipeline order. Only if no declared contract fails does it compare the final release with the canonical decision. This distinction is why a compiler failure remains attributed to the compiler when the database later contains it.

A mutant is \emph{killed} when this procedure emits a non-clean witness. It \emph{survives} when no modeled contract fails. A mutant is \emph{equivalent} when an injected change produces no policy-observable difference for the chosen principal, query, and state; it is \emph{invalid} (or stillborn) when it cannot form a valid test case. Clean controls contain no injected drift and become false positives if a witness is emitted.

\paragraph{Attribution versus root cause.}
The reported location is the first observed violated transition, not a source-code root cause. We test this attribution interventionally on compound cases. Repairing the attributed surface must move or eliminate the first violation (sufficiency), while repairing a later, non-attributed surface must leave the earlier attribution intact (necessity). A teeth test that forces a constant wrong attribution fails these checks. This is stronger than matching a fixture label, but it remains relative to the simulator and distinct-surface mutation model.

\paragraph{Witness reduction.}
The reducer tries one edit at a time: remove a semantic-query dimension, remove a filter, or reset a non-default limit. It accepts an edit only after replay preserves the witness class, first failing surface, containment layer, release decision, the localized failed contract, and any semantic-difference or database-containment fact needed by the original witness. The search stops after 32 attempts or a fixed point. Standard support suites reach 1-minimality under this move set, with median semantic-IR reductions of roughly 2\% for seeded cases and 6\% for generated cases. One-minimal means no single remaining edit in this three-operation neighborhood preserves the witness; it does not mean globally shortest JSON, SQL, policy, or database state. Because generated inputs are already narrow, this is diagnostic compaction, not general delta debugging.

\section{Retargeting \policystrata}

\policystrata{} is a Python package and command-line tool. The deterministic paper path needs no LLM API key and runs with \texttt{scripts/reproduce-final.sh}. It writes materialized tasks, JSONL traces, summaries, witnesses, comparator results, ablations, and benchmark manifests. Optional PostgreSQL and ClickHouse services exercise database-policy fixtures but are not included in the 1720-case score.

\paragraph{What a new target must provide.}
Retargeting is a specification and adapter task, not automatic source-code analysis. The minimum domain package contains:

\begin{enumerate}
  \item a canonical policy mapping real principals or roles to allowed metrics, dimensions, time ranges, row and cost limits, and tenant scopes;
  \item a surface contract naming the policy-pipeline versions and the obligations each available surface accepts and emits;
  \item representative semantic queries or imported traces with principal identity, SQL or semantic IR, and release outcome;
  \item adapters or configuration that connect real names, tenant columns, semantic-model fields, database roles, and release boundaries to those declarations.
\end{enumerate}

Database-backed checking additionally needs a sanitized schema and seed state, or a read-only connection to a representative fixture. Real-DB assertions can constrain row counts, columns, allowed values, and forbidden values. A target without a semantic model, RLS, or release observer may omit that input, but the resulting report must mark the surface unobserved rather than infer conformance.

\paragraph{Specification effort.}
The effort depends on how much policy is already machine-readable. A stack with a dbt-style semantic model, explicit tenant columns, structured tool traces, and RLS DDL can transform much of the input mechanically. A stack whose policy lives in prose requires a human to choose a canonical interpretation before testing begins. The tool cannot infer that interpretation without changing the claim from conformance testing to policy discovery. The artifact therefore reports each input as native, mechanically transformed, or synthesized.

The four brownfield targets give a concrete lower bound on this work:

\begin{center}
\scriptsize
\setlength{\tabcolsep}{2pt}
\begin{tabular}{@{}L{0.21\columnwidth}r r L{0.43\columnwidth}@{}}
\toprule
Target & Script lines & Traces & Main manual boundary \\
\midrule
MetricFlow & 378 & 68 & synthetic identity; no native tenancy \\
Midday & 57 & 5 & ORM SQL transcribed; no semantic IR \\
WrenAI & 222 & 2 & one rule mapped; synthetic role and bypass \\
Cube & 389 & 4 & two missing-filter regressions synthesized \\
\bottomrule
\end{tabular}
\end{center}

``Script lines'' is checked-in transform-script length, not person-hours or an optimized adapter estimate. MetricFlow additionally scans 19 test files and selects 68 of 266 test documents; it skips 43 multi-metric cases, 33 cases outside the chosen model, and 122 macro-dependent cases rather than inventing unsupported IR. Midday concatenates 39 native migrations and transcribes five cited ORM queries. WrenAI maps one required session-property rule. Cube maps one accepted policy and two already-rejected policy fixtures. These measurements show that retargeting is feasible but not push-button: each target requires an explicit decision about what cannot be represented.

\paragraph{Freeze and provenance.}
\texttt{freeze-benchmark} hashes detector source, generator source, all mutation specifications, policy and surface YAML, suite input, materialized tasks, package version, and Git commit. A run verifies those hashes before producing evidence. This prevents silent post-freeze tuning, but it does not make a generated suite independent. The paper uses ``detector-frozen generated'' rather than ``blinded'' for such cases.

\paragraph{Trusted computing base.}
Reference interpreters and adapters are trusted. An 18-operator adapter mutation study makes this risk concrete: 13 mutations hide or weaken real findings, three invent findings, one fails loudly, and one leaves the gate result unchanged. Thus 16/18 tested adapter faults silently corrupt scan output. The current mitigations are schema checks, adapter tests, source mapping, provenance hashes, and explicit evidence levels; independent canary traces and cross-checked gate computation remain future hardening.

\paragraph{Artifact support commands.}
The scanner imports SQL and semantic traces, compares dbt-style semantic metadata, optionally executes read-only database checks, and emits gateable findings. The doctor command inventories which policy, trace, database, prompt/tool, release, and CI inputs are wired. Doctor output is a configuration audit, not proof that the application invokes every control correctly.

\section{Benchmark Construction}

\bench{} contains three synthetic domains. \texttt{support\_saas} models tenant-scoped support and revenue analytics; \texttt{finance\_saas} models firm, advisor, household, and sensitive-identifier access; \texttt{analytics\_clickhouse} models project-scoped product analytics, aggregate-only roles, cohort thresholds, timezones, and materialized-view lineage. They vary schema and policy vocabulary but remain author-controlled fixtures.

\paragraph{Mutation registry.}
The registry contains 22 deterministic operators: 12 compiler, five database, two release, and one each for manifest, grammar, and validator. Nine produce over-permissive behavior, eight semantic drift, three lowering violations, and two unsafe releases. Examples include a stale metric alias, forbidden dimension admitted by the grammar, sensitive-column omission, dropped or stale tenant keys, gross/net or unique/count skew, join fan-out, fiscal/calendar mismatch, stale RLS ownership, missing project filters, materialized-view lineage loss, and unsafe small-cohort or sampled release. These are fault hypotheses, not a census of production incidents.

\paragraph{How the 1720 cases are made.}
The total combines 170 hand-authored seeded cases and 1550 algorithmic cases:

\begin{center}
\scriptsize
\setlength{\tabcolsep}{2pt}
\begin{tabular}{@{}L{0.37\columnwidth}rrL{0.27\columnwidth}@{}}
\toprule
Suite & Cases & Clean & Construction \\
\midrule
support seeded & 50 & 0 & hand-authored matrix \\
support generated & 500 & 0 & seed 1729; frozen run \\
support \texttt{heldout\_v1} & 500 & 0 & seed 260626; post-freeze \\
finance seeded & 20 & 0 & hand-authored matrix \\
finance \texttt{heldout\_v1} & 250 & 0 & seed 260626; post-freeze \\
analytics seeded & 100 & 0 & hand-authored matrix \\
analytics generated & 300 & 0 & seed 260626; frozen run \\
clean controls & 80 & 80 & seed 260627; no drift \\
\midrule
Non-clean total & 1720 & 0 & 170 manual + 1550 generated \\
\bottomrule
\end{tabular}
\end{center}

Seeded YAML matrices specify an operator, principal, request template, typed query, surface skew, and repetition count. The loader expands each matrix row into static cases. For generated case $j$, the algorithm chooses operator $o_j=O[j\bmod |O|]$ from the domain-applicable registry; selects the first non-admin principal; uses a seeded generator to choose allowed or denied metrics, dimensions, and time ranges needed to activate $o_j$; appends \texttt{-gen} to the affected surface version; creates the typed query; and shuffles all materialized tasks with the same seed. It does not mutate arbitrary application source, generate arbitrary SQL text, or mine production traces.

Cycling is balanced within a generated suite, not across the combined benchmark. Support and finance use 14 applicable operators; analytics uses all 22. Once hand-authored matrices and unequal suite sizes are combined, per-operator counts range from 13 to 121. Appendix~\ref{app:operator-accounting} reports every operator's surface, class, and count so coverage density is auditable rather than implied by the total.

\paragraph{Representativeness and mutant accounting.}
The registry was authored to cover recurring drift shapes across identity, capability exposure, semantic compilation, database containment, and release. This rationale is tested in two separate ways rather than inferred from 1720/1720. First, 25 cited public faults were screened; 19 could be mapped without reversing their documented direction onto 12 existing operators, while six were excluded as out of scope. These fixtures model the resulting drift shape, not the original vulnerable code path. Second, a spec-blind authoring pass produced 42 cases from contract documents without detector access; \policystrata{} agreed on 39, while three exposed an unspecified cost-combination rule.

The standard generator deliberately emits non-equivalent, well-formed cases, so equivalent and invalid counts are both zero. The runner nevertheless reports killed, survived, equivalent, invalid, clean-control, and false-positive counts separately. This design avoids silently discarding future equivalent or stillborn mutants, but the current zero counts should not be read as evidence that equivalent-mutant selection is solved.

\paragraph{Frozen is not independent.}
All 1550 generated cases are verified against a freeze manifest in the final reproduction. Of these, the 750 \texttt{heldout\_v1} cases carry post-detector-freeze provenance; the remaining 800 use the ordinary generated suite name but are still hash-pinned for the final run. Every case shares the same operator taxonomy and generator logic. Freezing makes the exact evaluation reproducible and exposes post-hoc changes; it does not provide third-party authorship or statistical generalization.

\section{Evaluation Results}

We organize the evaluation by the question each evidence source can answer. Deterministic suites measure fault-model coverage; spec-blind, external-source, and historical-revision suites probe taxonomy dependence; brownfield and deployment-linked checks measure retargeting behavior; intervention and adapter studies test diagnostic assumptions.

\paragraph{RQ1: Does the checker cover its declared fault model?}
All 1720 non-clean cases emit a witness, and 80 no-drift controls emit none. The result includes 170 hand-authored cases and 1550 generated cases. Because expected behavior and generated cases share the mutation registry, 1720/1720 is a regression invariant over the implementation, not an estimate of field recall.

\begin{center}
\scriptsize
\setlength{\tabcolsep}{2pt}
\begin{tabular}{@{}L{0.38\columnwidth}rrrr@{}}
\toprule
Suite group & N & Killed & Surv. & FP \\
\midrule
support: seeded/generated/frozen & 1050 & 1050 & 0 & -- \\
finance: seeded/frozen & 270 & 270 & 0 & -- \\
analytics: seeded/generated & 400 & 400 & 0 & -- \\
clean controls & 80 & -- & -- & 0 \\
\bottomrule
\end{tabular}
\end{center}

\paragraph{RQ2: What do transition contracts catch?}
We retain simple point checks as diagnostic ablations, but do not present them as competitors. The stronger comparison is a conventional specification-derived test suite with six checks: tenant predicate presence, rejected forbidden metric, rejected forbidden dimension, row-limit enforcement, blocked release after canonical denial, and golden values for named headline metrics. It is implemented without access to the mutation list. A second comparator performs pairwise differential checks over adjacent surface decisions, following the model-versus-implementation testing pattern used by authorization engines.

\begin{center}
\small
\begin{tabular}{@{}L{0.58\columnwidth}rr@{}}
\toprule
Checker & Caught & Missed \\
\midrule
\policystrata{} responsibility contracts & 1720 & 0 \\
Conventional spec-derived tests & 1579 & 141 \\
Layered point-control stack & 1561 & 159 \\
Pairwise surface differential & 899 & 821 \\
\bottomrule
\end{tabular}
\end{center}

The conventional suite's 141 misses comprise 70 semantic drifts on metrics without golden values, 38 unsafe releases of canonically allowed queries, and 33 containment failures invisible in released values. The pairwise differential misses cases where allow/deny decisions agree while compiler semantics drift; 613 of its misses are compiler-localized semantic drift. All three comparators report 0/80 false positives on the standard clean controls. These results show an observability gap under the benchmark model, not superiority to Margrave, VeriEQL, or a production policy engine, which answer different questions and were not reimplemented.

\paragraph{RQ3: Does non-generator evidence expose misses?}
Yes. A spec-blind pass hand-authored three cases for each of 14 support-domain operators after reading the contract and operator descriptions but not detector, simulator, or generator source. The detector killed 39/42. All three misses target the same cost-expansion operator: the public contract lists component costs and a maximum, but does not define their composition, so the author could not construct a reliably over-budget case. Surface attribution matched 42/42, while six predeclared containment guesses were also wrong. This exercise was performed by a non-independent author in the same development environment and is only a proxy for external authorship.

A separate source ledger screened 25 cited public faults from PostgreSQL, Supabase, ClickHouse, Cube, Looker, dbt/MetricFlow, Metabase, Superset, Hasura, and LangChain. Nineteen were mapped to 12 existing operators and all 19 fixtures were killed; six were dropped because their direction or mechanism was outside the model. Exact Git-object replay over three historical BetterOff fixes also verified that the relevant source contract was absent before and present after each fix. Two missing-RLS revisions map to \texttt{app\_deny\_missing\_db\_policy}; an export-audit omission is reproducible but outside v1. This replays historical source states, not running vulnerable services, and does not show prevention. Section~\ref{sec:external-taxonomy} adds the complementary scope check: measured against an independently authored data-agent vulnerability taxonomy, the v1 registry covers two of eight classes and one partially, and 401 of its own cases describe update-induced drift that an attack taxonomy has no reason to name.

\begin{center}
\scriptsize
\setlength{\tabcolsep}{2pt}
\begin{tabular}{@{}L{0.25\columnwidth}rL{0.28\columnwidth}L{0.31\columnwidth}@{}}
\toprule
Evidence & N & Independence & What it supports \\
\midrule
Generated mutants & 1550 & shared generator/taxonomy & regression coverage only \\
Seeded mutants & 170 & author-written & named examples only \\
Spec-blind & 42 & detector source hidden; same author & contract clarity and misses \\
Public-fault map & 25/19 & external reports; mapped model & taxonomy grounding \\
External taxonomy & 8 & independently authored classes & registry scope versus a second opinion \\
Brownfield SQL & 74 & external repositories & adapter behavior and precision cases \\
Executed real RLS & 20 & upstream policy text; local bridge & real policy enforcing and failing \\
MetricFlow freeze & 68 & upstream tests; local adapter & source-case reproducibility and adapter limits \\
Historical revisions & 3/3 & exact BetterOff Git objects & source-contract reproduction \\
Production HTTP & 33/36 & exact deployed revision & live public and denial boundaries \\
\bottomrule
\end{tabular}
\end{center}

No row in this table is an independently operated production study. The separate LASM cross-check maps the same 1720 cases onto a 116-paper architectural taxonomy; it is a scope analysis, not another case suite. The 68 MetricFlow rows are the frozen subset of the 74 brownfield traces, not an additional denominator. The table prevents calling upstream-authored cases an externally operated blind suite, treating historical source probes as vulnerable-service execution, or treating denial probes as field recall.

\paragraph{RQ4: Can the scanner be retargeted?}
Static scans targeted MetricFlow, Midday, WrenAI, and Cube. Inputs were labeled native, mechanically transformed, or synthesized. Across 74 traces carrying real SQL content, the initial pass produced one content-level false positive (1.4\%) and no newly discovered upstream defect. It did catch Cube's two already-known broken access-policy fixtures while leaving the valid fixture clean. A fresh MetricFlow checkout at Git object \texttt{45dce78641bb} reproduced all 68 transformed traces byte-for-byte. Its upstream-authored requests and expected SQL are more independent than generated cases, but the policy bridge remains author-written; 68 fuzz mutations survived because the synthetic bridge role permits every dimension. The five adapter gaps that pass exposed have since been fixed, which is the more useful outcome than the null discovery result: MetricFlow's adapter-attributable warnings fall from 27 to four, and each survivor is a finding the adapter should report. Running a scanner against code its authors did not write mostly measures the scanner.

Those scans are static, which leaves open whether the checks would notice a real policy that stopped working. One target answers that directly. Midday commits real PostgreSQL row-level security, so we load all 20 \texttt{CREATE POLICY} statements across the six policy-bearing tables in its frozen migrations verbatim into PostgreSQL~18.4 and execute them. Its migrations assume a Supabase base schema they do not commit, so a labeled bridge supplies the roles, \texttt{auth.uid()}, the team-lookup function, and the referenced base tables; checks connect as a non-owner, non-superuser role so that RLS applies rather than being silently bypassed. With Midday's policies intact, thirteen real-database checks pass and the gate is clean. Replacing one real predicate with \texttt{USING (true)}---the \texttt{db\_rls\_old\_ownership\_field} shape applied to real policy text---fails three RLS checks and one state assertion, including a read that is now available unauthenticated, while the five tables whose policies were untouched stay clean. This is real committed policy code enforcing and then failing to enforce, with a negative control; it is not a Midday defect, since its committed predicate is correct, and not its deployed Supabase runtime, which we have not observed.

\paragraph{RQ5: What does a deployment-linked pilot establish?}
Vercel reported BetterOff production deployment \texttt{dpl\_5MQsJfsc} ready at Git object \texttt{3663f1e475eb}; the primary checkout and adapter matched that revision. Thirty-three of 36 read-only HTTP probes passed and none failed. They cover health, authentication metadata, unauthenticated app/API/export denial, internal-route denial, invalid-webhook rejection, and public pages. Three authenticated reads were skipped because no isolated production session or API token exists; valid webhook probes were excluded because they can enqueue work. On the same revision, doctor found 33 tools, three roles, six runtime events, and no partial wiring; six SQL traces plus four database checks passed against disposable synthetic PostgreSQL. No customer row was read and no production state was mutated. Thus the pilot establishes deployment identity, adapter applicability, and live denial boundaries---not authenticated cross-tenant behavior, customer-data safety, or effectiveness on unknown faults.

\paragraph{RQ6: Are attributions causal?}
Counterfactual repair was evaluated on 120 compound cases per domain. In all 360 cases, repairing the attributed surface removed or moved that first violation, while repairing another surface left it in place. A forced constant-wrong localizer fails this test. The result validates first-transition attribution within the distinct-surface simulator; it does not establish source-code root-cause accuracy.

The bounded reducer reaches 1-minimality under its three edit types on the standard support suites, but reductions are small because generated plans already contain one dimension and few optional fields. Finally, the 18-case trusted-adapter mutation study finds 16 silent corruptions. This negative result sets a practical limit on every scan result: adapter correctness is part of the claim.

\paragraph{Cost.}
The frozen artifact run completes in about four seconds on the reported Apple laptop with a prepared environment, producing 1800 traces and 1720 witness files in a 17 MB directory. This is a reproducibility measurement on one machine, not a scaling law.

\section{Evaluation Protocol}

Every study is deterministic and reports raw counts because the standard suites are enumerated artifacts rather than samples from a defined production distribution. We do not attach confidence intervals to 1720 operator-generated cases.

\paragraph{Scoring.}
A non-clean case is killed only when the detector emits a non-clean witness. Survived, equivalent, and invalid cases remain in the denominator accounting rather than being discarded. Clean controls are scored separately. Expected class and expected surface are fixture labels; agreement with them is a regression check, not independent validation. Counterfactual repair is used when making the stronger attribution claim.

\paragraph{Information access.}
All comparators receive the same traces but only the fields implied by their design. The conventional suite reads specification-derived properties and selected golden values. The pairwise differential reads surface decisions but not mutation labels, first-transition labels, or witnesses. Point-control ablations inspect only their named surface. \policystrata{} alone evaluates the declared transition contracts; this asymmetry is the feature being tested and is stated rather than hidden.

\paragraph{Provenance separation.}
Reports keep five evidence kinds separate: hand-authored synthetic fixtures, operator-generated fixtures, detector-frozen generated fixtures, spec-blind hand-authored fixtures, and public-fault reconstructions. Imported real traces and real-database observations are scanner evidence, not automatically benchmark evidence. A freeze manifest proves source identity, not independence.

\paragraph{Reproducibility.}
The final script runs the suites, verifies freeze manifests, emits comparator and ablation reports, and records package and Git versions. Individual scripts reproduce the public-fault, spec-blind, counterfactual, minimization, scalability, clean-control, and adapter studies. No reported study needs an LLM API key. The model-reachability harness has not been run with a model and is excluded from all results.

\section{Artifact and Data Availability}

The implementation, paper source, benchmark inputs, study scripts, and per-study limitations are available at \url{https://github.com/raintree-technology/policystrata}. The repository also provides a Python package, command reference, evaluator guide, expected outputs, freeze manifests, optional database services, and scanner examples. Synthetic data are checked in; public-fault mappings retain their source ledger; brownfield adapters retain field-level native/transformed/synthesized provenance. The paper PDF is built from \texttt{paper/main.tex}, making the package repository the canonical editable source rather than the website.

\newpage

\section{What Remains Unestablished}

The maintainer-operated deployment study verifies source-to-deployment binding and selected unauthenticated denial boundaries, but its private identifiers are omitted from the public artifact. It does not exercise an authenticated tenant-to-tenant path because no isolated synthetic production principal is configured. Its deeper semantic, SQL, RLS, and release checks use the deployed source revision with synthetic rows. It therefore does not show public reproducibility, customer-data safety, model choice, retries, caches, concurrent rollout, unknown-fault recall, or independently operated effectiveness. External teams have not operated the study or authored a \policystrata-blind suite.

One external target's real policy text does now execute rather than being read statically, and it fails the right checks when a real predicate is weakened. That covers all 20 policies across the six policy-bearing tables in Midday's frozen migrations under a bridge we wrote, so it establishes that the containment checks bind to real policy code, not that they generalize to Supabase-managed schema or a deployed runtime.

The next evidence step is therefore an isolated synthetic principal for the three skipped authenticated probes, followed by a suite written by an independent author after detector freeze and an independently operated pilot. The historical replay harness should also advance from exact source-contract probes to executable vulnerable-service replays where dependencies and data fixtures permit.

Technical extensions also remain: same-surface compound mutations, mechanized contract proofs, verifier-backed SQL equivalence, wider witness reduction, sequential and history-aware release, adaptive differencing, writes and side effects, and model-mediated reachability under fixed prompt, model, retry, and refusal settings. Results for these extensions should remain separate from the v1 read-only benchmark.

\section{Threats to Validity}

\paragraph{Construct validity.}
The operator taxonomy may omit real fault shapes, and the simulator may make represented faults easier to detect. The spec-blind cost misses show that even an implemented operator can be underspecified. Finite database states can expose a semantic difference but cannot prove query equivalence.

\paragraph{Internal validity.}
Generated inputs, expected labels, and detector behavior share code and concepts. Freeze manifests prevent silent changes but not common design errors. Reference interpreters reduce compiler-oracle coupling, while the adapter mutation study shows that bad normalization can still hide or invent findings.

\paragraph{External validity.}
All scored benchmark domains are synthetic. Public-fault cases map outcomes, while the three BetterOff historical cases probe exact pre/post-fix source contracts rather than execute vulnerable services. Brownfield work covers four open-source repositories and uses synthesized bridge values. The one executed-policy result covers all 20 policies across the six policy-bearing tables in Midday's frozen migrations, but its Supabase runtime is our reconstruction rather than the deployed one, so it shows real policy text enforcing under a bridge we wrote. The BetterOff pilot is bound to production but its deep checks use synthetic data, and its authenticated probes were skipped. No customer data or independently operated deployment was evaluated. Two external-taxonomy cross-checks bound the registry against independently authored vulnerability and architectural vocabularies, but both mappings are our judgement and neither taxonomy is a field distribution.

\paragraph{Conclusion validity.}
Counts are exact for checked-in deterministic suites but do not estimate a production fault distribution. Zero false positives on generated controls is weak precision evidence because those controls share the simulator. The one-of-74 brownfield figure and 33 passing live probes are more independent but too small, adapter-dependent, and boundary-focused to generalize.

\section{Conclusion}

\policystrata{} treats a data agent as a policy pipeline whose representations are related but not equal. Explicit transition contracts make it possible to distinguish an upstream lowering error, downstream containment, and final release rather than collapsing them into one pass/fail answer. The artifact covers its declared 22-operator model, exposes gaps missed by specification-derived tests, and emits intervention-checked first-transition witnesses. Spec-blind misses, null brownfield discovery, and silent adapter mutations also show where the evidence stops. The result is a reproducible method and retargeting contract for policy-drift regression testing, not a claim of production recall or a replacement for enforcement.

\appendix

\section{Contract Formalization}
\label{app:formalization}

The formalization fixes terminology for the implemented contracts rather than providing a soundness theorem for arbitrary data-agent deployments.

Let $p$ be a principal and $c$ be runtime context containing request metadata, session state, and deployment versions. Let $X_i$ be the object language at surface $i$, and let $\gamma_i(x)$ interpret object $x \in X_i$ as the concrete operations, data dependencies, or observations it may induce. Let $A_P(p,a,c)$ mean operation or dependency $a$ is authorized by $P^t_{\mathsf{auth}}$. Let $\mathsf{Expose}_i$, $\mathsf{Accept}_i$, and $\mathsf{Release}$ describe actual implementation behavior. Let $E_i(p,x,c)$ be the declared support envelope for surface $i$, and let $L_i:X_i\rightharpoonup X_{i+1}$ be a partial lowering.

\paragraph{Exposure soundness.}
\[
\mathsf{Expose}_i(p,x,c) \Rightarrow \forall a \in \gamma_i(x), A_P(p,a,c)
\]

\paragraph{Declared completeness.}
\[
E_i(p,x,c) \wedge \forall a \in \gamma_i(x), A_P(p,a,c) \Rightarrow \mathsf{Accept}_i(p,x,c)
\]

\paragraph{Authorization-preserving lowering.}
\[
\mathsf{Accept}_i(p,x,c) \wedge L_i(x)=y
\Rightarrow
\forall a \in \gamma_{i+1}(y), A_P(p,a,c)
\]

\paragraph{Semantic translation validation.}
\[
\exists D: \mathsf{Obs}(\mathsf{Exec}(L_i(x),D)) \not\approx \mathsf{SpecEval}(x,D)
\]

\paragraph{Release conformance.}
\[
\mathsf{Release}(p,y,\ell,c) \Rightarrow R_P(p,y,\ell,c)
\]

\paragraph{Checked detector invariant.}
Let $\mathsf{Fail}_i(T)$ mean that the implemented contract decision for surface $i$ rejects trace $T$, and let $\mathsf{UnauthorizedRelease}(T)$ mean the release decision allows an observation rejected by the canonical policy. The implementation checks
\[
\mathsf{Witness}(T)\Rightarrow
\left(\bigvee_i\mathsf{Fail}_i(T)\right)
\vee \mathsf{UnauthorizedRelease}(T).
\]
This invariant is executable and tested, not proved for arbitrary adapters or implementations.

\paragraph{Why the implementation satisfies the implication.}
The detector first searches the ordered contract-decision map. If it returns surface $i$, then by construction $\mathsf{Fail}_i(T)$ is true, so the right side holds. If no surface contract fails, the only branch that emits a non-clean witness checks that release is allowed while the canonical decision denies it. Otherwise the detector returns \textsc{clean}. Classification into over-permissive, lowering, semantic, or release subclasses happens only after one of these two guards. This is a proof sketch over the detector control flow, while the property and finite exhaustive tests check that trace construction supplies the premises as expected.

\paragraph{Completeness boundary.}
The 22 operators cover three lowering violations at the compiler, nine over-permissive faults across manifest, grammar, validator, compiler, and database, eight compiler semantic drifts, and two unsafe releases. Completeness is asserted only for well-formed cases instantiated by these operators. No global converse to the checked invariant is claimed.

\section{Operator Accounting}
\label{app:operator-accounting}

The table lists the full v1 read-only mutation registry and its frequency in the 1720 non-clean cases. Counts combine all seeded and generated suites; unequal domain sizes explain why operators are not globally balanced.

\begin{center}
\scriptsize
\setlength{\tabcolsep}{2pt}
\begin{tabular}{@{}L{0.49\columnwidth}llr@{}}
\toprule
Operator family & Surface & Class & N \\
\midrule
stale metric alias & manifest & over-perm. & 121 \\
forbidden dimension admitted & grammar & over-perm. & 121 \\
sensitive column omitted & validator & over-perm. & 121 \\
dropped tenant predicate & compiler & lowering & 111 \\
old tenant key & compiler & lowering & 109 \\
tenant/account id swap & compiler & lowering & 106 \\
old RLS ownership field & database & over-perm. & 109 \\
gross/net metric skew & compiler & semantic & 111 \\
fan-out join & compiler & semantic & 109 \\
removed distinct & compiler & semantic & 106 \\
inner join drops rows & compiler & semantic & 104 \\
fiscal/calendar mismatch & compiler & semantic & 108 \\
cost expansion ignored & compiler & over-perm. & 104 \\
app deny absent in DB & database & over-perm. & 106 \\
missing project filter & database & over-perm. & 25 \\
read-only assumption broken & database & over-perm. & 23 \\
small cohort released & release & unsafe & 25 \\
materialized-view lineage loss & compiler & semantic & 25 \\
timezone bucket skew & compiler & semantic & 25 \\
unique/count skew & compiler & semantic & 25 \\
sampled result released & release & unsafe & 13 \\
distributed-table policy gap & database & over-perm. & 13 \\
\midrule
Total & & & 1720 \\
\bottomrule
\end{tabular}
\end{center}


\balance
\begingroup
\footnotesize
\bibliographystyle{plainnat}
\bibliography{references}
\endgroup

\end{document}
